\chapter{RESULTS INTERPRETATION}

\section{Analysis of Behavioral Finance Integration}
The primary success of this study lies in the dynamic performance gain when incorporating \textbf{FinBERT sentiment analysis} into the feature set. This confirms core tenets of \textbf{Behavioral Finance} \citep{shiller2003}, which suggests that market prices are driven more by collective psychology and news-driven sentiment than by immediate fundamental shifts. Sentiment polarity, fetched via Finnhub, exhibits a notable predictive capacity during high-volatility events across MSFT, AMZN, and GOOGL.

\section{Theoeretical Comparison Across Asset Classes}
\begin{itemize}
    \item \textbf{High-Growth Equities:} Theoretical sensitivity to information asymmetry is captured by the model, as evidenced by improved recall when sentiment scores are present. This aligns with the "Investor Sentiment" theory \citep{baker2007}, which posits that stocks with difficult-to-value future earnings are most prone to sentiment shifts.
    \item \textbf{Commodities (Gold):} Exhibits defensive behavior that contrasts with equity volatility. Sentiment signals tied to systemic risk (hedging) proved to be more predictive than traditional technical indicators during market downturns.
    \item \textbf{Digital Assets (Bitcoin):} High volatility leads to wider confidence intervals. However, the model exploits the irrational exuberance and "herd behavior" often observed in crypto markets, where sentiment polarity serves as a leading indicator of momentum.
\end{itemize}

\section{What the Numbers Mean: The "Reality Check"}
The metrics obtained (Accuracy 46--61\% and OOS Sharpe Ratios typically $< 1.0$, except for Gold) represent a transparent and scientifically rigorous evaluation after adjusting for **market friction** and **evaluation bias**. 

A directional accuracy of ~50\% for major tech equities suggests that at a daily frequency, price movements in these highly liquid assets closely resemble a **Random Walk**. However, the XGBoost model identifies localized momentum in Gold (GC=F) with an OOS Sharpe of 1.49, suggesting that its defensive volatility patterns may be more predictable using the current feature set.

We consider the model "Trustworthy" not because it predicts high returns, but because its evaluation protocol—using strictly isolated test sets and transaction costs—provides a conservative and reliable estimate of actual performance. This contrasts with "frictionless" models that falsely report high Sharpe ratios due to in-sample leakage and zero-cost assumptions.

\section{Limitations and Stochastic Volatility}
While the system is robust, the results emphasize the stochastic nature of finance, where "Black Swan" events or high-frequency news bursts can override historical signals. The non-stationarity identified in our ADF tests confirms that the model's predictive alpha is not static and requires continuous walk-forward retraining to mitigate model decay.
